\begin{claim}{D4.1}{Proved}{Integral Picard twists, including their torsion}
For $M_n=\OO_X(n(S_1-S_2))$, define
$u(n)=\max\{d(n)+1,0\}$, $v(n)=\max\{-d(n)-1,0\}$ and
$\epsilon_n=0$ if $3\mid n$, $\epsilon_n=1$ otherwise. Then
\[
 f_*(T_X\otimes M_n)=\OO(d(n))\oplus\OO(d(n-1)),
\]
\[
 R^1f_*(T_X\otimes M_n)=
 \OO(d(n))\oplus\OO(d(n-1))\oplus\OO(d(n+2)),           \tag{D.12}
\]
\[
 R^2f_*(T_X\otimes M_n)=
 \OO(d(n+6))\oplus\C_{p_1}^{\epsilon_n}.
\]
In particular its cusp stalk is free of rank one. The four global
dimensions, for every integer $n$, are
\begin{align*}
 h^0&=u(n)+u(n-1),\\
 h^1&=v(n)+v(n-1)+u(n)+u(n-1)+u(n+2),\\
 h^2&=v(n)+v(n-1)+v(n+2)+u(n+6)+\epsilon_n,\\
 h^3&=v(n+6).                                          \tag{D.13}
\end{align*}
\end{claim}
\begin{proof}
The defining roots give $\bar f^*R=\OO_X(S_1-S_2)$, so the projection
formula applies to (D.11), with its actual coherent torsion summand.
Tensoring (D.8) changes its generator character from $-1$ to $n-1$.
There is exactly one invariant coefficient among exponents $0,1$ if
$n-1$ is $0$ or $1$ modulo three, and none if it is $2$. Multiplication
by the coarse parameter $t=s^3$ is zero on this module. Its invariant
direct image is therefore exactly the displayed skyscraper, of
length $\epsilon_n$, not a higher-order thickening. The two other
special places contribute no torsion. Formula (D.1) proves (D.12),
and the two-column Leray sequence proves (D.13). The code evaluates
these formulas for $-240\le n\le240$ using two independent integer
cohomology counts; their alternating sums are all one. That last
check is supplementary to the sheaf proof.
\end{proof}
